\chapter[Sermon 93. Here Is Set Forth a Goodly Picture, a Description or Figure of the Blessed Seat, and of the Heavenly Life and Glory Everlasting.]{Here Is Set Forth a Goodly Picture, a Description or Figure of the Blessed Seat, and of the Heavenly Life and Glory Everlasting.}
\label{sermon:093}
\markright{Sermon 93. Here Is Set Forth a Goodly Picture, a Description or ...}

And one of the seven angels, who had the seven vials full of the seven last plagues, came to me and spoke, saying, “Come here, and I will show you the bride, the Lamb’s wife.” And he carried me away in spirit to a great and high mountain, and he showed me the great city, holy Jerusalem, descending out of heaven from God, having the glory of God. And her shining was like to a stone most precious, even a jasper clear as crystal; and had great and high walls, and twelve gates, and at the gates twelve angels; and names written, which are the names of the twelve tribes of the children of Israel. Three gates on the east side, three gates on the north side, three gates toward the south, and three gates on the west side. And the wall of the city had twelve foundations, and in them the names of the Lamb’s twelve apostles.

John returns to the description of the heavenly city, which at the beginning of this chapter he had begun. He has included certain matters that are fitting, concerning the certain hope of the faithful; after which he seems afterward to unlock and open Heaven, that the godly, with the eyes of faith, might as it were look therein, and clearly see what is the hope and glory of saints to come. For under the image of a most beautiful city, he sets forth a picture or description most evident of the blessed dwelling place, or palace and city of God, or of the everlasting realm and triumphant church. We shall not here invent and fabricate things earthly and physical, but spiritual and celestial. For the Spirit of God intends that, by means of temporal things, our minds ascend to the eternal, and that by temporal things, we grasp what is more excellent. Therefore all things are figured, with amplifications, hyperboles, and full of other figures. We shall therefore imagine far greater things: as we are accustomed to do, when we read or hear such things as our Lord taught under the parables of weddings and feasts.

\index{Rome}And first, it is declared to us who is the revealer of this godly and wonderful vision: that is to say, who is the opener of the mysteries, truly an Angel of God, and the very same who previously in chapter 17 to the same John said, “Come, I will show you the condemnation of the great whore,” and so forth. For it is the same God who punishes the ungodly, and gives rewards to the godly, and announces to men his righteous judgments by his ministers. Moreover, we see the sins to be most certain, and partly also accomplished, which he showed before concerning the judgment of Rome. Would not anyone conclude that what he now utters and shows concerning the everlasting glory of the faithful is also most certain? Gathering some of the things which he will show him, he sets them before and exhorts him to follow, saying: “Come, I will show you the bride, the wife of the Lamb.” Of her there has been spoken before on many occasions. He signifies the congregation of Saints, coupled by faith to our Savior Christ. And not only does he show to John (and in the same, to us all) the spouse, but also reveals her glory from God. The meaning therefore is this: “Come, I will show you what shall be the glory of the church of Christ in the life to come, what shall be the state of the everlasting life.” Certainly, he speaks also very many things of the church, but chiefly of her glory in the world to come.

\index{Patmos}\index{Arethas of Caesarea}\index{Ezekiel (Book of)}Then he briefly touches on the manner of revelation. For he adds, “And he took me up in spirit onto a great and high mountain.” Therefore, just as in the former visions he was carried away in spirit, his body remaining in Patmos, and as we have read and pointed out before, that such visions and raptures happened to Ezekiel, so now he says that he is carried away in spirit, and in mind to have seen the things the angel showed. Wherefore, if we will read or hear these things to any profit, we must lift up our minds, and be carried up in our spirit, and think that all these things must be understood spiritually. Arethas: rightly, he says, the heavenly life and conversation of the saints was shown to him on the mountain. For with them there is nothing earthly, low, or abject, but all things lofty and high. This he. Certainly, when the Lord Christ wished to exhibit to his disciples a certain taste and foretaste of the glory to come, he carried them up onto a mountain and was transfigured before them, which Matthew affirms in chapter 17, to have happened to Peter, James, and John.

And now he turns to the vision itself, and generally and briefly describes or foreshadows the blessed seat and glory of the life to come. Afterward, he amplifies this more largely, particularly, and as it were by parts, and so enlarged and beautified he sets it forth as if it were to be seen by the godly.

\index{Aquinas, Thomas}He calls the heavenly country, and the habitation of saints, the great City. For it is the city of the great king, and in it shall dwell an innumerable number of the blessed, and of angels thousands without end, and shall have the fruition of great glory. Neither is there any fear that the place should not suffice so great a host of men and spirits, or that it shall be too crowded. Great is the city of God, which is truly able to receive all good men abundantly. In the Gospel of John the Lord says: “In my Father’s house are many mansions,” and so forth, John 14:2. The same place is called holy Jerusalem. For just as no defilement shall be seen there, so shall no unclean person appear. Of the heavenly Jerusalem is spoken before. Thomas Aquinas says: She is said to have descended from heaven, for that whatever goodness the holy church has, she acknowledges that she has received it from the grace of God. But of this matter I have spoken in the last Sermon.

And the city of God, that is to say heaven, has the seats of God and the blessed, the glory of God, that is to say the divine majesty and brightness, and whatever great thing so ever the mind of man can think or imagine, or in all things the unspeakable excellence of God, such as neither the eye has seen, nor the ear has heard, nor has ascended into the heart of man, 1 Corinthians 2:9. These things he has touched upon briefly and generally hitherto.

And consequently he declares by particulars and at large that celestial glory and blessed seats are present. For whatever things are simple, whatever are in cities commendable, the same are plainly found in this our city most excellent, such as its size, strength, majesty, security, excellence, beauty, pleasantness, and abundance of things. These things I say, and all other like, wonderfully excel in the city of our God and in our fathers’ house. And whereas these things are set forth and amplified in this way most liberally, yet it seems nothing at all is said, if a man considers the unspeakable majesty of the celestial glory. But all these things are alleged by the Lord through John to this end, truly, that we should be taken with the desire of so worthy a life, and should think in our tribulations and troubles that the afflictions of this present world are nothing, being compared with so excellent and sovereign glory; finally, that all are mad, who begin to doubt of the eternal hope of the faithful. Very many things of this sort are also read in Ezekiel in chapter 40 and following. We will touch every part of this treatise, using nevertheless a succinct brevity, lest we should be tedious to any man. And truly he touches the principal and most commendable things of cities, and in them shows that the city of God excels.

In cities and houses, the greatest commendation is when all things are bright and clear, for darkness is horrible and unpleasant. Therefore, an excellent light is declared to be in the city or house of the Lord. A parable is added, showing the excellence of this light. It is like a most precious stone, such as a jasper, as it is commonly called, or a chrysolite, or some other stone most bright. And John himself adds more, as it were a jasper stone like a crystal. This is a new way of speaking, but it has a marvelous grace if we understand it rightly. For a jasper is green, and a crystal is bright. He seems therefore to say that this celestial brightness is continually green and never withers; that is to say, the heavenly light is everlasting, and in itself, after a sort, wearing green, and in growing green, it shines brightly and rejoices all heavenly dwellers. For afterward follows: “For the glory of God has lightened her, and the Lamb is her light.” This brightness and most joyful light the Lord promises in sundry places in the gospel of John, and the whole blessed life, of this not the least part, is commonly called blessed light, or everlasting light, or light of heaven. It seems to have been prefigured in the golden candlestick of the tabernacle, and so forth. For if it were not difficult for our Lord God to give to precious stones wonderful colors and brightness; if he illuminates this world full of wicked men with most goodly lights, the Sun, Moon, and Stars, what light may we think to have in heaven, where no man shall dwell but the best, and of God most dearly beloved? Of this light much mention is made with Isaiah and in the psalmist.

Walls in cities are most notable and excellent when they are high, thick, and strong, able to withstand all force from enemies, and defend the citizens from all harm, keeping them in peace and security. Therefore, strong walls are both great and strong, and also high and impregnable. This signifies that the protection of the saints in heaven will be most safe and sure, so that the saints will be in perfect security, and exempt from all fear. No one will trouble or take away their joys, as the Lord affirmed in John 16. For there will be perpetual security and gladness in heaven, most perfect and everlasting.

Moreover, gates are placed in the walls, by which people may enter the city. Therefore, in the wall of the great country there shall be twelve gates, that is to say, a most ample entrance into eternal life shall be open on every side. And we believe that there is no other way to heaven, no other port or gate, or any other door or postern to remain, than the only and sole Christ Jesus our Lord, as he himself has taught in John 10 and 14. But because he has appointed angels, prophets, and apostles also, as porters of heaven, to whom he has committed the keys of the kingdom of heaven, and these do bring the chosen and let them into the heavenly country, many gates are truly said to have been and to be. And for a further declaration, it is added that in every gate was an angel, namely twelve. And we have heard in the beginning of this book that angels are God’s ministers and pastors of churches, sent by God for the salvation of people, that is, that they might bring them by the word of truth and holy ministry, through faith, into life everlasting. Moreover, we read how the soul of poor Lazarus, dying, was carried by angels into the bosom of Abraham. Why, then, should we marvel that angels stand at the gates? For by the true and only gate, Christ, they bring the faithful into the heavenly country.

\index{Hebrews, Epistle to the}And again, for a further declaration is appended, and in the gates were names written, which are the names of the twelve tribes of the children of Israel. For the Lord would signify that he used the industry of the patriarchs and prophets of all tribes in opening heaven to me; and again, that all the elect of all tribes belong to the fellowship of blessedness. We shall see, therefore, in heaven, the patriarchs and prophets, and all the saints who, before the coming of Christ, are written in the registers of the heavenly realm, just as the apostles also saw Moses and Elijah talking with Christ on the mount. Wherefore, not without great cause, the Apostle to the Hebrews wrote: “You have come to Mount Zion, and to the city of the living God, the heavenly Jerusalem, and to the multitude of angels, and to the assembly of the firstborn who are written in heaven.” And the rest which is read in the twelfth chapter.

And he also addresses the matter of the gates. He assigns three to each part of the sky, and does not do this without reason. For our Savior himself says in the Gospel that they will come from the east and from the west, and will rest with Abraham, Isaac, and Jacob, in the kingdom of heaven. Arethas investigates this mystery more closely, and believes that no one will enter through these gates unless they acknowledge the eternal Trinity of God and understand the mystery of the cross of Christ. He says how the twelve tribes are divided by the Trinity, according to the fourfold figure of the world, \&c. Aquinas also says: whoever are saved, they are justified by the faith of the holy Trinity, proclaimed throughout the four quadrants of the world by the preaching of the Apostles.

Now he also shows that the foundations of this city are most sure and immovable. For the wall of the City, he says, has twelve foundations. Concerning the foundation of the church and our salvation, David spoke expressly in the Psalms. Isaiah also in chapter 28. Our Lord and Savior spoke in sundry places in the Gospel. Peter moreover in the Acts and his first epistle: likewise the Apostle Paul, who said, “No other foundation can be laid than that which is laid, which is Christ Jesus,” 1 Corinthians 3. How then are there laid here twelve foundations? Doubtless Christ remains one and a sure foundation. However, inasmuch as in placing and revealing him, the Lord used the ministry of the twelve Apostles, for this cause the city is said to have twelve foundations. Not that the Apostles are in deed the foundations of the church and our salvation, but in this that Christ, that true foundation, was by the twelve Apostles made known to the faithful, and as it were laid under, whereupon the believers have built themselves by the Apostles’ faith. Whereupon he says purposefully, “and in those twelve are the names of the Lamb’s twelve Apostles.” For the gospel also, which is both in very deed and unchangeably Jesus Christ’s alone, is called of John, Matthew, Mark, and Luke, of Peter and Paul, because it has been preached by them. And we understand hereby not only the church which was before the coming of Christ, of Patriarchs and Prophets, being received into heaven to rejoice in God, but also the Apostolic church, I mean that all men in the whole world who have believed the Apostolic doctrine shall live with all the Saints in that heavenly country; all of which we shall both see, and with them also shall glorify God forevermore.

\index{Primasius}Primasius, Bishop of Utica, does not greatly disagree with this our explanation, explaining how the Apostles are called foundations. For thus he has written: since we know that the church has one only foundation, that is to say Christ, we ought not to be disturbed that here he says she has twelve. For in Christ the Apostles have deserved to be the foundations of the church; of whom the Apostle says, “another foundation can not be laid, besides that which is laid, Christ Jesus.” In him also the Apostles are said to be light, as he says to them, “You are the light of the world,” where Christ alone is the true light, which illuminates every man coming into this world. Christ therefore is the illuminating light, and they are the lights illuminated. And after a few words, the same author: “Here it behooves to acknowledge the twelve Apostles to be foundations, yet called in the only foundation, Christ Jesus.” Hereunto pertains also that he has not concealed the name of the Lamb. The Apostles therefore are foundations, but in one foundation, Jesus Christ. And Christ alone, without the Apostles, is rightly called the foundation; but the Apostles without Christ, could by no means be called the foundations of the church. These are the words of Primasius. Which Arethas, Bishop of Caesarea, declares more briefly and plainly, and says: they are indeed called foundations, for that they have laid the foundations of the Christian faith; and gates, for that by them—that is to say, by their preaching—there may be found a way to bring people to the Christian faith. Thus much he says. Doubtless the Apostle Paul in 2 Ephesians calls Christ the foundation of Apostles and Prophets, which verily in preaching they have laid, and to which they have leaned, and by which also they are saved. To him be glory.
